\section{Propagation of automorphy and descent}\label{sec:automorphy}

We now fix a parameter $t\in U(K)$ supplied by
Proposition~\ref{res:specialization}, taking the local open sets of
Proposition~\ref{loc:opens}. Write $B=B_t$. The four systems
$V_\eta,V_\chi,V_\psi,V_1$ have the congruences of
Corollary~\ref{sys:congruences}, the local properties of
Corollary~\ref{loc:characters}, and the residual properties proved in
Section~\ref{sec:residual}. We first explain why automorphy at one
coefficient prime supplies the information needed at the next one.

\subsection{Coefficient primes and local automorphic conditions}

\begin{lemma}\label{aut:transport}
Suppose a weight-zero regular algebraic cuspidal representation $\pi$ of
$\GL_2(\A_F)$ realizes one member of the compatible system $V_\theta$.
Then it realizes every member, using the compatible coefficient embeddings
fixed in Section~\ref{sec:systems}.
\end{lemma}
\begin{proof}
At all places outside a finite set and away from the coefficient prime,
the characteristic polynomial of Frobenius on $V_\theta$ is independent
of that prime, as an algebraic polynomial with the fixed embeddings of the
character values. The good-place characteristic polynomials of
$r_\iota(\pi)$ have the same compatibility. Equality at the original prime
therefore identifies the polynomials at every other prime. Both
representations are semisimple: for $V_\theta$ this follows from Faltings'
theorem as in Proposition~\ref{sys:properties}, and the automorphic Galois
representations are taken semisimple. Chebotarev density and the
Brauer--Nesbitt theorem give the asserted isomorphism.
\end{proof}

The next lift also requires information about the local automorphic
components at the new coefficient prime. We recover it by applying
Varma's comparison at a different coefficient prime, where it applies
directly to the local Weil representation.

\begin{lemma}\label{aut:automorphic-local}
Suppose that a weight-zero regular algebraic cuspidal representation
$\pi$ of $\PGL_2(\A_F)$ realizes $V_\theta$. Let $l\in\{3,p\}$ and
$v\mid l$, and write $J_t=\Jac(C_n(t))$.
\begin{enumerate}[(i)]
\item If $J_t$ has potentially good reduction at $v$, then
$\rec_{F_v}(\pi_v)$ has zero monodromy.
\item If $J_t$ acquires semistable reduction of toric rank $n$ at $v$,
then $\pi_v$ is a
quadratic or trivial twist of the Steinberg representation, and $\pi$ is
$\iota_l$-ordinary of weight zero at $v$.
\end{enumerate}
\end{lemma}
\begin{proof}
Choose a coefficient prime $a$ different from the residue characteristic
$l$. Lemma~\ref{aut:transport} identifies $r_{\iota_a}(\pi)$ with
$V_{\theta,a}$. Varma's local-global compatibility theorem identifies the
underlying semisimplified Weil representations of its local
Weil--Deligne parameter and $\rec^T_{F_v}(\pi_v)$
\cite[Theorem~1]{Varma}. We use this statement about the Weil
representations first; equality of the monodromy operators is not assumed.

In the potentially good case, pass to an extension where the Jacobian has
good reduction. All Frobenius eigenvalues on its $H^1$, and hence on the
character summand, have complex absolute value the square root of the
residue cardinality. Taking a positive power of Frobenius shows that the
two eigenvalues before this extension also have the same absolute value.
A rank-two Weil--Deligne representation with nonzero monodromy has
Frobenius eigenvalues differing by the norm character. They cannot have
the same absolute value. Thus the automorphic parameter has zero
monodromy. The Tate twist in \eqref{eq:tate-normalization} does not change
this conclusion.

In the totally toric case, Corollary~\ref{loc:characters} gives
nonzero monodromy on $V_{\theta,a}$. Its semisimplified Weil representation
is therefore a sum of two characters differing by the norm. For this
Weil representation, the two possibilities under the local Langlands
correspondence for $\GL_2$ are the special parameter with nonzero
monodromy and the zero-monodromy parameter of a one-dimensional local
representation. The latter is not generic. Since every local component
of a cuspidal automorphic representation of $\GL_2$ is generic, $\pi_v$
is a twist of Steinberg. Its central character is trivial, so the twisting
character has square one.

The local calculation in the proof of
\cite[Lemma~6.1.3(4)]{CN}, which invokes the proof of
\cite[Lemma~5.6]{Geraghty}, shows that a quadratic twist of Steinberg is
$\iota_l$-ordinary of weight zero. This calculation uses only the stated
local representation and trivial central character. In particular, it
also applies when the quadratic character is ramified: the ordinary
condition is imposed at a sufficiently deep level, and a finite-order
twist multiplies the relevant normalized Hecke eigenvalues by units.
\end{proof}

\subsection{The three lifts}

\begin{proposition}\label{aut:over-F}
The elliptic curve $B/F$ is modular. More precisely, $H^1(B)$ is realized
by a cuspidal weight-zero regular algebraic representation of
$\PGL_2(\A_F)$.
\end{proposition}
\begin{proof}
The order-three system $V_\eta$ is $H^1(A)$ for the elliptic curve
$A/F$ of Proposition~\ref{sys:elliptic}. It is non-CM by
Proposition~\ref{res:specialization}. Its mod-$5$ representation is
decomposed generic and absolutely irreducible over $F(\zeta_5)$ by the
$(m,l)=(3,5)$ part of that proposition. Since $\zeta_5\notin F$,
Theorem~\ref{aut:seed} supplies a cuspidal weight-zero representation of
$\PGL_2(\A_F)$ realizing $V_\eta$.

Apply Theorem~\ref{aut:lift} in the following order:
\begin{center}
\begin{tabular}{@{}cccc@{}}
\toprule
Automorphic input & Coefficient prime & System lifted & Residual input\\
\midrule
$V_\eta$ & $p$ & $V_\chi$ & $(m,l)=(3,p)$\\
$V_\chi$ & $3$ & $V_\psi$ & $(m,l)=(p,3)$\\
$V_\psi$ & $p$ & $V_1$ & $H^1(B)\bmod p$\\
\bottomrule
\end{tabular}
\end{center}
In each row Corollary~\ref{sys:congruences} identifies the residual
representations, and Proposition~\ref{res:specialization} gives decomposed
genericity and absolute irreducibility over the cyclotomic extension.
Proposition~\ref{sys:properties} supplies the determinant, finite
ramification, potential semistability, and labeled Hodge--Tate weights.
Lemma~\ref{aut:transport} makes the previously obtained automorphy
available at the coefficient prime of the next row.

It remains to check the local hypotheses at that prime. At a place of
potentially good ordinary type, both systems in the row are potentially
crystalline and potentially ordinary of weight zero. At a place of
potentially good supersingular type, they are potentially crystalline
and neither is potentially ordinary of weight zero. These assertions
follow from Corollary~\ref{loc:characters}. In either case
Lemma~\ref{aut:automorphic-local}(i) gives zero monodromy in the input
automorphic parameter. At a place of potentially totally toric type,
neither Galois representation is potentially crystalline, and
Lemma~\ref{aut:automorphic-local}(ii) gives automorphic ordinarity.
Thus the relevant alternative of Theorem~\ref{aut:lift} holds at every
place above the coefficient prime.

The theorem therefore applies in all three rows. Its conclusion preserves
cuspidality, weight zero, and trivial central character, so each row
provides the input required for the next. The last row gives automorphy of
$V_1=H^1(B)$.
\end{proof}

\subsection{Descending to the original field}

We recall why automorphy over the cyclotomic extension suffices here.
The character adjustment in the following argument is important: it
identifies the descended Galois representation with the prescribed one.

\begin{lemma}\label{aut:descent}
Let $L'/L$ be a cyclic extension of prime degree between CM fields. Let
$r:G_L\to\GL_2(\overline\Q_l)$ be a continuous representation whose
restriction to $G_{L'}$ is absolutely irreducible. Suppose that this
restriction is realized by a cuspidal weight-zero regular algebraic
representation $\Pi$ of $\GL_2(\A_{L'})$. Then $r$ is realized by a
cuspidal weight-zero regular algebraic representation of $\GL_2(\A_L)$.
\end{lemma}
\begin{proof}
For $\sigma\in\Gal(L'/L)$, the representations $\Pi^\sigma$ and $\Pi$
have the same good-place Hecke polynomials, because the associated
Galois representation extends to $G_L$. Strong multiplicity one makes
$\Pi$ Galois invariant. The cyclic base-change theorem therefore gives
a cuspidal automorphic representation $\pi$ of $\GL_2(\A_L)$ whose
base change is $\Pi$~\cite{LanglandsBC,ArthurClozel}.
Every archimedean place is complex, so its local extension is trivial.
Consequently $\pi$ is regular algebraic of weight zero.

Good-place compatibility with base change and Chebotarev imply
\[
 \bigl(r_\iota(\pi)|_{G_{L'}}\bigr)^{\semisimple}
 \simeq r|_{G_{L'}}.
\]
The right-hand side is absolutely irreducible. The restriction on the
left is therefore already irreducible and isomorphic to it, without
semisimplification. After choosing this isomorphism, the two
representations agree on $G_{L'}$. For $g\in G_L$, their matrices at $g$
differ by a scalar, since they induce the same conjugation on the
absolutely irreducible representation of $G_{L'}$. Schur's lemma shows
that these scalars form a character
$\delta:\Gal(L'/L)\to\overline\Q_l^{\times}$ and
\[
 r_\iota(\pi)\simeq r\otimes\delta.
\]
The finite-order Hecke character corresponding to $\delta^{-1}$ removes
this discrepancy. Twisting preserves cuspidality and weight zero, and
the compatibility with Galois twisting follows at good places and hence
everywhere by Chebotarev.
\end{proof}

Since $F/K$ is abelian, choose a tower from $K$ to $F$ with cyclic
prime-degree steps. Each intermediate field is abelian over $\Q$ and
contains the imaginary quadratic field $K$, hence is CM. The
representation $H^1(B,\overline\Q_p)$ remains absolutely irreducible
over $F$: its residual representation is already absolutely irreducible
there. Apply Lemma~\ref{aut:descent} down the tower to
Proposition~\ref{aut:over-F}. This gives automorphy of $B/K$. Finally,
$B$ is the quadratic twist of $E$ by $h(t)$, so twisting the automorphic
representation by the same quadratic Hecke character gives a cuspidal
weight-zero representation $\pi_E$ realizing $H^1(E)$.

\subsection{Local parameters and the CM case}

For non-CM $E$, Theorem~\ref{aut:elliptic-compat} shows that $\pi_E$
has trivial central character, realizes $H^1(E)$ for every coefficient
prime, and has the required Frobenius-semisimplified Weil--Deligne
parameter at each finite place, by choosing a coefficient prime different
from the residue characteristic. Thus every finite local $L$-factor
matches. At infinity a cuspidal representation is generic. The
archimedean classification with weight-zero infinitesimal character
therefore gives exactly the parameter \eqref{eq:infinity-parameter};
the finite-dimensional alternative is not generic. Equations
\eqref{eq:tate-normalization} and \eqref{eq:L-shift} give the stated
normalization of the local and global $L$-functions.

Suppose now that $E$ has complex multiplication by an imaginary quadratic
field $M$. Its endomorphisms are defined over $KM$. One way to see this
is to use their faithful action on invariant differentials: the
eigenvalues lie in $M$, and are fixed by $G_{KM}$. The main theorem of
complex multiplication expresses the two one-dimensional constituents
of $H^1(E)|_{G_{KM}}$ as the realizations of algebraic Hecke
characters, with the infinity types prescribed by the CM
type~\cite{ShimuraCM}.

If $M\not\subset K$, Galois exchanges the two characters and $H^1(E)$
is induced from either one over the quadratic extension $KM/K$.
Automorphic induction gives the required representation of
$\GL_2(\A_K)$. If $M=K$, use the isobaric sum of the two characters.
In either case normalize the resulting representation by
$|\det|^{1/2}$ so that its Tate-normalized parameter realizes $H^1(E)$.
Complex multiplication and local class field theory, together with local
compatibility of quadratic automorphic induction in the first case,
give the finite parameters and the prescribed infinity types. These are
the CM alternatives allowed in Theorem~\ref{thm:main}. Their central
character is trivial: the determinant of $H^1(E)$ is $\eps_l^{-1}$,
so undoing the Tate normalization gives determinant one for the local
automorphic parameters, and hence trivial central character.

In both the CM and non-CM cases, the finite factors also agree if one uses the
Weil--Deligne representation at the residue-characteristic coefficient
prime. For good reduction this follows from crystalline comparison. For
bad potentially good reduction, the Weil--Deligne inertia action is finite
and nontrivial, and monodromy is zero. The inertia action has determinant
one. Its invariant subspace cannot have dimension one, since a
semisimple two-dimensional determinant-one representation with a
nonzero fixed vector would be trivial. It therefore has no invariants, giving
local factor one. At the residue-characteristic coefficient prime,
nontriviality of inertia follows from the crystalline good-reduction
criterion~\cite{ColemanIovita}. Finally, in the potentially multiplicative
case the Tate-curve description, allowing a quadratic twist, gives the
same special Weil--Deligne parameter at all coefficient primes. This
completes the proof of Theorem~\ref{thm:main}.
